\documentclass[11pt]{article}
\usepackage{amsmath,amssymb,amsthm,bm}
\usepackage{geometry}
\geometry{margin=1in}
\usepackage{setspace}
\usepackage{enumitem}
\usepackage{booktabs}
\usepackage{mathtools}
\usepackage{hyperref}
\usepackage{algorithm}
\usepackage{algpseudocode}

\title{Joint Multilayered LSIRM with Item-Only Factor-Mixture Clustering\\
(Model V8)}
\date{April 29, 2026}

\begin{document}
\maketitle

\section{Overview}

Model V8 couples the 5-layered hierarchical LSIRM of V6 with an item-only
factor-mixture clustering (FMC) layer in a single Bayesian hierarchical model.
The structure is summarized as follows.
\begin{itemize}[leftmargin=1.5em]
    \item \textbf{LSIRM block.} Identical to V6.
        Five mixed-type response layers (binary, robust continuous, count with overdispersion, two ordinal layers) share a hierarchical respondent-position structure $(z_i,\,a_i^{(1)},\dots,a_i^{(5)})$.
        Item positions $b_j^{(\ell)}\in\mathbb{R}^d$ live in the same latent space.
        Latent positions are constrained to a unit disk centered at the origin.
    \item \textbf{Clustering block.} An item-only finite mixture in factor-score
        space, applied to the respondent-relation profile of each item induced
        by the current LSIRM state.
    \item \textbf{Coupling.} The current implementation is strict one-way
        (LSIRM~$\to$~FMC). The LSIRM acceptance ratios are
        byte-identical to V6; the clustering block conditions on the LSIRM
        latent positions but does not feed back into LSIRM updates.
\end{itemize}

The remainder of the document fixes notation, formalizes the joint model,
specifies hyperparameters (with explicit rationale for the choices that
matter for cluster recovery), and documents diagnostic failure modes that
were observed in fits to the MIDUS data along with the corrective changes
adopted.

\section{Notation}

\subsection*{LSIRM block}

\begin{itemize}[leftmargin=1.5em]
    \item $n$: number of respondents, indexed $i=1,\dots,n$.
    \item $L=5$: number of data layers.
    \item $P_\ell$: number of items in layer $\ell$; $P=\sum_{\ell=1}^L P_\ell$ is the total number of items, indexed $j=1,\dots,P$. Let $\ell(j)$ denote the layer of item $j$.
    \item $d$: dimension of the latent interaction space (fixed at $d=2$).
    \item $z_i \in \mathbb{R}^d$: \emph{global} respondent latent position; appears only as the hierarchical mean of the local positions and never enters any likelihood directly.
    \item $a_i^{(\ell)} \in \mathbb{R}^d$: \emph{local} (layer-specific) respondent latent position in layer $\ell$.
    \item $b_j^{(\ell)} \in \mathbb{R}^d$: latent position of item $j$ in its corresponding layer.
    \item $\sigma_1^2 > 0$: global--local coupling variance, controlling how tightly $a_i^{(\ell)}$ clusters around $z_i$.
    \item $\alpha_i^{(\ell)}$: respondent effect for respondent $i$ in layer $\ell$.
    \item $\beta_j^{(\ell)}$: item difficulty; for ordinal layers, replaced by GRM thresholds $\beta_{j,1}^{(\ell)}>\cdots>\beta_{j,K_\ell-1}^{(\ell)}$.
    \item $\gamma_\ell>0$: layer-specific distance weight.
    \item $\kappa>0$: count-layer overdispersion.
\end{itemize}

For brevity, layer-specific variances $\sigma_{\alpha\ell}^2,\tau_{\beta\ell}^2,\sigma_0^2$ and the robust-continuous mixture weights $\lambda_{ij}^{(2)}$ follow the V6 specification verbatim.

\subsection*{Clustering block (FMC)}

\begin{itemize}[leftmargin=1.5em]
    \item $r$: latent factor dimension.
    \item $K^\star$: overfitted upper bound on the number of item clusters.
    \item $\eta_j\in\mathbb{R}^r$: factor score of item $j$.
    \item $c_j\in\{1,\dots,K^\star\}$: cluster label of item $j$.
    \item $K_+^{(t)}=|\{c_j^{(t)}:j=1,\dots,P\}|$: number of \emph{occupied} components at MCMC iteration $t$.
    \item $\Lambda\in\mathbb{R}^{n\times r}$: shared factor-loading matrix; row $\lambda_i^\top$ is respondent~$i$'s loading.
    \item $\delta=(\delta_1,\dots,\delta_n)^\top$: per-respondent baseline.
    \item $\Psi=\mathrm{diag}(\psi_1,\dots,\psi_n)$: per-respondent idiosyncratic variances.
    \item $\rho=(\rho_1,\dots,\rho_{K^\star})$: mixture weights.
    \item $\mu_l\in\mathbb{R}^r$, $\Sigma_l\in\mathbb{R}^{r\times r}_{\succ 0}$: mixture component parameters.
    \item $\sigma_\delta^2$: variance of the respondent baselines.
\end{itemize}

\section{LSIRM block (recap of V6)}

The hierarchical respondent-position prior is
\begin{align}
z_i &\overset{iid}{\sim} N_d(0_d,I_d),\\
a_i^{(\ell)}\mid z_i,\sigma_1^2 &\overset{ind}{\sim} N_d(z_i,\sigma_1^2 I_d),
\quad \ell=1,\dots,L,\\
\sigma_1^2 &\sim IG(a_{\sigma_g},b_{\sigma_g}).
\end{align}

Per-layer linear predictors and likelihoods are unchanged from V6 and are reproduced compactly:
\begin{align*}
\eta_{ij}^{(1)} &= \alpha_i^{(1)}-\beta_j^{(1)}-\gamma_1\|a_i^{(1)}-b_j^{(1)}\|, & Y_{ij}^{(1)}&\sim\mathrm{Bern}(\mathrm{logit}^{-1}\eta_{ij}^{(1)}),\\
\mu_{ij}^{(2)} &= \alpha_i^{(2)}-\beta_j^{(2)}-\gamma_2\|a_i^{(2)}-b_j^{(2)}\|, & Y_{ij}^{(2)}&\sim t_{\nu_2}(\mu_{ij}^{(2)},\sigma_0^2)\;\text{(via normal--gamma mixture)},\\
\mu_{ij}^{(3)} &= \alpha_i^{(3)}-\beta_j^{(3)}-\gamma_3\|a_i^{(3)}-b_j^{(3)}\|, & Y_{ij}^{(3)}&\sim\mathrm{NB}(\kappa,\exp \mu_{ij}^{(3)}),\\
\mu_{ij}^{(4)} &= \alpha_i^{(4)}-\gamma_4\|a_i^{(4)}-b_j^{(4)}\|, & Y_{ij}^{(4)}&\sim \text{GRM}(\mu_{ij}^{(4)};\beta_{j,1}^{(4)},\dots,\beta_{j,K_1-1}^{(4)}),\\
\mu_{ij}^{(5)} &= \alpha_i^{(5)}-\gamma_5\|a_i^{(5)}-b_j^{(5)}\|, & Y_{ij}^{(5)}&\sim \text{GRM}(\mu_{ij}^{(5)};\beta_{j,1}^{(5)},\dots,\beta_{j,K_2-1}^{(5)}).
\end{align*}

The full LSIRM likelihood is
\[
p(Y\mid\Theta_{\mathrm{LSIRM}})
=\prod_{\ell=1}^L\prod_{i=1}^n\prod_{j:\,\ell(j)=\ell}
p_\ell\bigl(Y_{ij}^{(\ell)}\mid\theta_{ij}^{(\ell)}\bigr).
\]

\section{Item relation profile}

For each MCMC iteration, the current LSIRM configuration induces, for each item~$j$, a respondent-relation profile vector
\[
x_j(\Theta_{\mathrm{LSIRM}})
=
(x_{1j},\dots,x_{nj})^\top \in \mathbb{R}^n,
\qquad
x_{ij} = \log \|a_i^{(\ell(j))}-b_j^{(\ell(j))}\|.
\]
Although the LSIRM uses layer-specific local positions $a_i^{(\ell)}$, the FMC sees a single $n$-vector per item; the layer index simply selects which $a_i^{(\ell)}$ is used in the distance computation.
For notational simplicity we drop the layer superscript on $a_i$ in the FMC equations below.

\paragraph{Row centering (mandatory preprocessing).}
At the start of every MCMC iteration the FMC input is row-centered:
\[
\widetilde{x}_{ij}=x_{ij}-\frac{1}{P}\sum_{j'=1}^P x_{ij'},\qquad
\sum_{j=1}^P \widetilde x_{ij}=0\;\forall i.
\]
Without this step, the per-respondent baseline $\delta_i$ absorbs the respondent-mean log-distance (i.e.\ $\sigma_\delta^2\sim\mathrm{Var}(x_{\log})$), leaving the $\Lambda\eta_j$ block with negligible signal and forcing the FMC to collapse to $K_+=1$.
Row-centering is therefore not optional: it is a structural requirement for the cluster signal to survive.
We use $\widetilde x_j$ as the FMC input throughout but write $x_j$ from now on for brevity.

\section{Factor-mixture clustering block}

Conditional on the LSIRM state, the FMC layer specifies
\begin{align}
x_j \mid \delta,\Lambda,\eta_j,\Psi
&\sim N_n\bigl(\delta+\Lambda\eta_j,\;\Psi\bigr),\label{eq:fmc-likelihood}\\
\eta_j \mid c_j=l,\mu_l,\Sigma_l &\sim N_r(\mu_l,\Sigma_l),\\
c_j \mid \rho &\sim \mathrm{Categorical}(\rho),\\
\rho &\sim \mathrm{Dirichlet}(e_0,\dots,e_0).
\end{align}

\subsection{Priors}

\begin{align}
\delta_i\mid\sigma_\delta^2 &\sim N(0,\sigma_\delta^2), &
\sigma_\delta^2 &\sim IG(a_\delta,b_\delta),\\
\lambda_i &\sim N_r(0,\tau_\lambda^2 I_r),\\
\psi_i &\sim IG(a_\psi,b_\psi),\\
\mu_l &\sim N_r(m_0,V_0), &
\Sigma_l &\sim IW(\nu_0,S_0).
\end{align}

\subsection{Hyperparameter choices and rationale}\label{sec:hyper-rationale}

The defaults below are used in all V8 fits to the MIDUS data ($n=221$, $P=68$, $r=5$, $K^\star=10$, $e_0=0.1$).

\begin{table}[h]
\centering
\begin{tabular}{lll}
\toprule
Parameter & Value & Rationale \\
\midrule
$e_0$           & $0.1$       & Overfitted Dirichlet (Rousseau--Mengersen): $e_0<d_l/2$ where $d_l=r+r(r+1)/2=20$ for $r=5$, so any $e_0<10$ asymptotically empties redundant components. \\
$m_0$           & $0_r$       & Center prior on cluster centers at the origin. \\
$V_0$           & $9I_r$      & Loose prior on $\mu_l$, large mobility. \\
$\nu_0$         & $r+10=15$   & Informative IW prior; peaks $\Sigma_l$ near $S_0$. \\
$S_0$           & $0.05\,I_r$ & Tight prior on within-cluster covariance, so distinct $\mu_l$'s correspond to distinct clusters. \\
$\tau_\lambda^2$ & $0.25$     & Small; loosening $\Lambda$ further reverts the $\Lambda\to 0$ collapse. \\
$(a_\delta,b_\delta)$ & $(2,1)$ & $E[\sigma_\delta^2]=1$, weak prior.\\
$\bm{(a_\psi,b_\psi)}$ & $\bm{(5,0.2)}$ & \emph{Tightened from $(2,1)$.} See Section~\ref{sec:fixA}. Forces $\psi_i$ near $0.05$ so that residual variance must flow through $\Lambda\eta_j$ rather than into $\Psi$. \\
\bottomrule
\end{tabular}
\caption{Hyperparameter defaults for the FMC block in V8.}
\end{table}

\section{Joint hierarchical model and coupling}

The joint posterior kernel is
\begin{align*}
p(\Theta_{\mathrm{LSIRM}},\Theta_{\mathrm{clust}}\mid Y)
&\propto
p(Y\mid\Theta_{\mathrm{LSIRM}})\\
&\quad\times
\prod_{j=1}^P
N_n\!\bigl(x_j(\Theta_{\mathrm{LSIRM}});\delta+\Lambda\eta_j,\Psi\bigr)
N_r(\eta_j;\mu_{c_j},\Sigma_{c_j})\rho_{c_j}\\
&\quad\times
p(\Theta_{\mathrm{LSIRM}})\,p(\rho)\,
p(\delta\mid\sigma_\delta^2)\,p(\sigma_\delta^2)\,p(\Lambda)\,p(\Psi)
\prod_{l=1}^{K^\star}p(\mu_l)p(\Sigma_l),
\end{align*}
where
\[
\Theta_{\mathrm{clust}}=\{\eta_j,c_j,\rho,\delta,\Lambda,\Psi,\mu_l,\Sigma_l,\sigma_\delta^2\}.
\]

\paragraph{One-way coupling.}
In V8 the Metropolis--Hastings acceptance ratios for $a_i$ and $b_j$ omit the FMC term~$N_n(x_j;\delta+\Lambda\eta_j,\Psi)$.
That is, the LSIRM block updates as if the FMC did not exist, and the FMC sees the LSIRM through $x_j(\Theta_{\mathrm{LSIRM}})$ but does not influence it.
This sacrifices full posterior coherence for two-way coupling, but keeps the V6 LSIRM block byte-identical and avoids the implementation cost of recomputing the cluster contribution at every $a_i,b_j$ proposal.
A two-way version is sketched in Section~\ref{sec:future} as future work.

\section{Initialization (PCA-based)}\label{sec:init}

The FMC has a degenerate fixed point at $(\Lambda,\eta)=(0,0)$ that random small initial values cannot escape: with $\eta\approx 0$ the data favors $\Lambda\approx 0$, which in turn favors $\eta\approx 0$ via the mixture prior.
We break this chicken-and-egg using a PCA initialization:
\begin{enumerate}[leftmargin=2em,label=(\arabic*)]
    \item Stack the column-standardized observed responses across layers into an $n\times P$ matrix $X$ (NA-imputed at $0$).
    \item Compute PCA on $X^\top$ (items as rows) and retain the first $r$ scores as $\eta_j^{(0)}$.
    \item Run $K^\star$-means on $\eta^{(0)}$ to obtain initial cluster labels $c_j^{(0)}$ and centers $\mu_l^{(0)}$.
    \item Initialize $\Lambda^{(0)}\sim N(0,0.01\,I)$; the first $\lambda_i$ Gibbs sweep then redraws $\Lambda$ conditional on the structured $\eta^{(0)}$.
\end{enumerate}

\paragraph{FMC warm-up.}
For $t<T_{\mathrm{warm}}$ (default $T_{\mathrm{warm}}=\max(1000,\mathrm{burnin}/4)$), the FMC parameters are not updated; only LSIRM positions evolve.
This lets $a_i,b_j$ converge before the FMC starts conditioning on $x_j(\Theta_{\mathrm{LSIRM}})$, so that the FMC is not anchored to noisy initial distances.

\section{Posterior updates}\label{sec:posterior-updates}

LSIRM updates are unchanged from V6.
We list only the FMC-block conditionals.
Throughout, $\mathcal{J}_l=\{j:c_j=l\}$ and $n_l=|\mathcal{J}_l|$.

\paragraph{(1) $\rho$.}
$\rho\mid c\sim\mathrm{Dirichlet}(e_0+n_1,\dots,e_0+n_{K^\star})$.

\paragraph{(2) $\eta_j$.}
$\eta_j\mid\cdot\sim N_r(m_{\eta_j},V_{\eta_j})$ with
\[
V_{\eta_j}=\bigl(\Lambda^\top\Psi^{-1}\Lambda+\Sigma_{c_j}^{-1}\bigr)^{-1},\qquad
m_{\eta_j}=V_{\eta_j}\bigl(\Lambda^\top\Psi^{-1}(x_j-\delta)+\Sigma_{c_j}^{-1}\mu_{c_j}\bigr).
\]

\paragraph{(3) $c_j$.}
$P(c_j=l\mid\cdot)\propto\rho_l\,\phi_r(\eta_j;\mu_l,\Sigma_l)$, $l=1,\dots,K^\star$.

\paragraph{(4) $\mu_l$.}
$\mu_l\mid\cdot\sim N_r(m_{\mu_l},V_{\mu_l})$ with
$V_{\mu_l}=(V_0^{-1}+n_l\Sigma_l^{-1})^{-1}$ and
$m_{\mu_l}=V_{\mu_l}\bigl(V_0^{-1}m_0+\Sigma_l^{-1}\sum_{j\in\mathcal{J}_l}\eta_j\bigr)$.

\paragraph{(5) $\Sigma_l$.}
$\Sigma_l\mid\cdot\sim IW\bigl(\nu_0+n_l,\;S_0+\sum_{j\in\mathcal{J}_l}(\eta_j-\mu_l)(\eta_j-\mu_l)^\top\bigr)$.

\paragraph{(6) $\delta_i$.}
$\delta_i\mid\cdot\sim N(m_{\delta_i},V_{\delta_i})$ with
$V_{\delta_i}=(P/\psi_i+1/\sigma_\delta^2)^{-1}$ and
$m_{\delta_i}=V_{\delta_i}\Bigl[\psi_i^{-1}\sum_{j=1}^P(x_{ij}-\lambda_i^\top\eta_j)\Bigr]$.

\paragraph{(7) $\lambda_i$.}
$\lambda_i\mid\cdot\sim N_r(m_{\lambda_i},V_{\lambda_i})$ with
$V_{\lambda_i}=\bigl(\tau_\lambda^{-2}I_r+\psi_i^{-1}\sum_j\eta_j\eta_j^\top\bigr)^{-1}$ and
$m_{\lambda_i}=V_{\lambda_i}\Bigl[\psi_i^{-1}\sum_j\eta_j(x_{ij}-\delta_i)\Bigr]$.

\paragraph{(8) $\psi_i$.}
$\psi_i\mid\cdot\sim IG\bigl(a_\psi+P/2,\;b_\psi+\tfrac12\sum_j(x_{ij}-\delta_i-\lambda_i^\top\eta_j)^2\bigr)$.

\paragraph{(9) $\sigma_\delta^2$.}
$\sigma_\delta^2\mid\cdot\sim IG\bigl(a_\delta+n/2,\;b_\delta+\tfrac12\sum_i\delta_i^2\bigr)$.

\paragraph{LSIRM positions under one-way coupling.}
$a_i,b_j$ are updated by Metropolis--Hastings with the V6 acceptance ratios; the FMC factor $N_n(x_j;\delta+\Lambda\eta_j,\Psi)$ is \emph{not} included.

\section{One full MCMC sweep}

\begin{algorithm}[H]
\caption{One sweep of joint LSIRM~+~FMC (V8, one-way coupling)}
\begin{algorithmic}[1]
\State Update LSIRM non-position block ($\alpha,\beta,\gamma,\kappa,\sigma_0^2$, threshold parameters, $\sigma_1^2$, $z_i$).
\State Update $\{a_i^{(\ell)}\}$ via MH (V6 acceptance, no FMC term).
\State Update $\{b_j^{(\ell)}\}$ via MH (V6 acceptance, no FMC term).
\State Form $x_{ij}=\log\|a_i^{(\ell(j))}-b_j^{(\ell(j))}\|$ and row-center to $\widetilde x_{ij}$.
\If{iter $\geq T_{\mathrm{warm}}$}
    \State Update $\rho$.
    \State Update $\{\eta_j\}_{j=1}^P$.
    \State Update $\{c_j\}_{j=1}^P$.
    \State Update $\{\mu_l,\Sigma_l\}_{l=1}^{K^\star}$.
    \State Update $\{\delta_i,\lambda_i,\psi_i\}_{i=1}^n$.
    \State Update $\sigma_\delta^2$.
    \State Update online co-cluster matrix $C_{jk}\leftarrow C_{jk}+\mathbb{1}\{c_j=c_k\}$.
\EndIf
\end{algorithmic}
\end{algorithm}

\section{Cluster inference}\label{sec:cluster-inference}

Per-iteration cluster labels $c_j^{(t)}$ are subject to label switching and are not directly interpretable.
Label-switching-invariant summaries are used:

\paragraph{Posterior similarity matrix (PSM).}
\[
C_{jk}=\Pr(c_j=c_k\mid Y)\;\approx\;\frac{1}{T}\sum_{t=1}^T \mathbb{1}\{c_j^{(t)}=c_k^{(t)}\}.
\]

\paragraph{$K_+$ trace.}
$K_+^{(t)}=|\{c_j^{(t)}\}|$. The point estimate of cluster count is
\(
\hat K=\max\bigl(2,\mathrm{round}(\mathrm{median}_t\,K_+^{(t)})\bigr).
\)

\paragraph{Final partition.}
A point partition is obtained by average-linkage hierarchical clustering on the dissimilarity $1-C$, cut at $\hat K$:
\[
\hat c_j = \mathrm{cutree}\bigl(\mathrm{hclust}(1-C,\,\text{method}=\text{average}),\;k=\hat K\bigr).
\]

\section{Diagnostics, failure modes, and remediations}\label{sec:diagnostics}

In MIDUS fits with the original psi prior $(a_\psi,b_\psi)=(2,1)$, the FMC consistently locks at $K_+=2$ across $e_0\in\{0.05,0.1,1\}$. We document the diagnosis here so future runs can be verified.

\subsection{Symptom: $K_+$ lock-in at $K=2$}\label{sec:symptom}

For the canonical run $(r=5,K^\star=10,e_0=0.1)$:
\begin{itemize}[leftmargin=1.5em]
    \item $\mathrm{median}\,K_+ = 2$, $\mathrm{sd}\,K_+ = 0$ across all 8000 saved iterations.
    \item Final partition is a clean $34$-vs-$34$ split.
    \item Sweeping $e_0$ does not change $\hat K$.
\end{itemize}
This persistent lock-in is not the well-known "$K_+=1$ trap" warned in the implementation notes (which occurs without row-centering); it is a deeper failure mode in which the mixture cannot resolve more than 2 components.

\subsection{Diagnosis: $\Lambda\to 0$ collapse}

The following diagnostics, computed on the posterior of the v8 fit, identify the cause.

\begin{itemize}[leftmargin=1.5em]
    \item \textbf{Loading norm.} $\|\Lambda_{\mathrm{PM}}\|_F=1.20$, while the prior expectation under $\tau_\lambda^2=0.25$ is $\sqrt{n r\tau_\lambda^2}\approx 16.6$. Lambda has collapsed to roughly $1/14$ of its prior scale.
    \item \textbf{Reconstruction.} $\Lambda\eta_j$ has range $[-0.007,0.008]$ and per-item standard deviation $0.001$; meanwhile the row-centered input $\widetilde x_j$ has standard deviation $\sim 0.5$. The factor block explains less than $2\%$ of the input variance.
    \item \textbf{Cluster separation vs.~spread.} For the two active components, $\|\mu_4-\mu_{10}\|=0.065$ but $\sqrt{\mathrm{tr}\,\Sigma_l}\approx 0.54$. Component Gaussians are essentially co-incident.
\end{itemize}

\paragraph{Mechanism.}
Latent positions $a_i,b_j$ are constrained to a unit disk, hence
$\|a_i-b_j\|\in[0,2]$, $\log\|a_i-b_j\|$ has very narrow dynamic range, and
$\widetilde x_{ij}$ has standard deviation of order $0.1$--$0.5$.
With the loose default psi prior $(a_\psi,b_\psi)=(2,1)$ (mean $1$, heavy tail), the posterior favors the simple decomposition
\[
\widetilde x_{ij}\approx \delta_i + \varepsilon_{ij},\qquad \varepsilon_{ij}\sim N(0,\psi_i),
\]
absorbing essentially all variance into $\psi_i$.
This drives $\Lambda\to 0$, which makes $\eta_j$ essentially data-free: each $\eta_j$ is then drawn from its mixture prior, and the mixture cannot discover structure that is not in the data.
The end state is two indistinguishable Gaussians in $\eta$-space and $K_+=2$.

\subsection{Fix A: tightening the $\psi$ prior}\label{sec:fixA}

We tighten the residual prior to $(a_\psi,b_\psi)=(5,0.2)$, giving
$E[\psi_i]\approx 0.05$ with concentration ($\mathrm{sd}(\psi_i)\approx 0.03$), versus the original mean $1$ with infinite variance.
The change forces the posterior to attribute the observed $0.5$-scale variance in $\widetilde x_{ij}$ to the structured term $\Lambda\eta_j$, restoring the factor block to a non-trivial role.
The other priors (in particular the small $S_0=0.05\,I_r$) do not need adjustment: once $\Lambda\eta_j$ carries variance, $S_0$ correctly biases $\Sigma_l$ toward small within-cluster spread, allowing $K_+>2$ to emerge.

\paragraph{Verification checklist after refit.}
\begin{enumerate}[leftmargin=2em]
    \item $\|\Lambda_{\mathrm{PM}}\|_F$ should rise toward the prior expectation~16.6 (target: $\geq 5$).
    \item $\mathrm{sd}(\widetilde x_j-\delta-\Lambda\eta_j)\approx \sqrt{\mathrm{mean}\,\psi_i}\approx 0.22$ initially, decreasing as $\Lambda\eta_j$ takes over.
    \item Active $\|\mu_l-\mu_m\|$ should exceed $\sqrt{\mathrm{tr}\,\Sigma_l}$ for at least one pair.
    \item $K_+^{(t)}$ should show non-degenerate variation (sd~$>0$).
\end{enumerate}

\subsection{Fix considerations not adopted}

\begin{itemize}[leftmargin=1.5em]
    \item \emph{Column standardization of $\widetilde x_j$}: rejected. Per-column rescaling can be absorbed by $\Lambda$ or $\eta_j$; it does not change the $\Lambda/\Psi$ trade-off and so does not address the mechanism.
    \item \emph{Removing row-centering}: rejected; collapses to $K_+=1$ as documented in the cpp implementation notes.
    \item \emph{Loosening $\tau_\lambda^2$}: rejected. Empirically reverts to a stable but data-disconnected $\Lambda\sim 0$ fixed point.
    \item \emph{Relaxing the unit-disk constraint on $a_i,b_j$}: a viable structural change but invasive (requires retuning all LSIRM proposals); kept as a fallback if Fix~A proves insufficient.
\end{itemize}

\section{Post-hoc clustering (sanity check and grid exploration)}\label{sec:posthoc}

Independently of the in-model PSM clustering of Section~\ref{sec:cluster-inference}, V8 stores the per-iteration factor scores $\eta^{(t)}\in\mathbb{R}^{P\times r}$ (`save_eta_full = TRUE`).
This enables post-hoc clustering on geometric distances to (i) verify that the in-model partition is supported by the factor-space geometry and (ii) probe alternative cluster counts on a $K=2,\dots,8$ grid.

\subsection{Distance in respondent-profile space}

Items live in $\mathbb{R}^r$ via $\eta_j$, but the substantive distance of interest is between their reconstructed respondent profiles $\hat x_j=\delta+\Lambda\eta_j\in\mathbb{R}^n$.
Since $\delta$ cancels in differences,
\[
\|\hat x_j-\hat x_k\|^2 = \|\Lambda(\eta_j-\eta_k)\|^2 = (\eta_j-\eta_k)^\top M\,(\eta_j-\eta_k),
\quad M=\Lambda^\top\Lambda.
\]
With $R$ the Cholesky factor of $M$ ($R^\top R=M$), defining $\tilde\eta_j=R\eta_j$ gives the clean transform
\[
\|\hat x_j-\hat x_k\|=\|\tilde\eta_j-\tilde\eta_k\|.
\]
Because per-iteration $\Lambda^{(t)}$ is not stored (the FMC saves $\Lambda_{\mathrm{PM}}$ only), $R$ is computed as $\mathrm{chol}(\Lambda_{\mathrm{PM}}^\top\Lambda_{\mathrm{PM}})$ and treated as a plug-in.

\subsection{Two clustering inputs}

\paragraph{(P1) Posterior-mean factor scores in respondent metric.}
\[
\eta_{\mathrm{PM},j} \;=\; R\,\bar\eta_j,
\quad
\bar\eta_j=\frac{1}{T}\sum_{t=1}^T \eta_j^{(t)}.
\]
Used as input to $K$-means and hierarchical clustering (Euclidean).

\paragraph{(P2) Posterior-mean pairwise distance.}
\[
\bar D_{jk} \;=\; \frac{1}{T}\sum_{t=1}^T \bigl\|R(\eta_j^{(t)}-\eta_k^{(t)})\bigr\|.
\]
Used as input to PAM and hierarchical clustering (precomputed dissimilarity).
By Jensen's inequality $\bar D_{jk}\geq \|\eta_{\mathrm{PM},j}-\eta_{\mathrm{PM},k}\|$, with equality only when $\eta_j-\eta_k$ is constant across iterations.
$\bar D$ thus preserves posterior uncertainty about pairwise positions that point estimation discards.

\subsection{Cluster-quality grid}

For each input we run hierarchical clustering with four linkages (ward.D2, average, complete, single) and $K$-means or PAM, sweeping $K=2,\dots,8$ and selecting $K$ by the average silhouette width, with elbow and gap statistics as cross-checks.
The resulting partitions are compared via the adjusted Rand index against the in-model PSM partition $\hat c$ and against pre-existing SVD-based factor-analysis partitions.

\subsection{Empirical findings on the canonical pre-fix run}

For the $(a_\psi,b_\psi)=(2,1)$ run:
\begin{itemize}[leftmargin=1.5em]
    \item Both inputs and all algorithms select $K=2$ at silhouette-optimal $K$ (with one exception: ward.D2 on $\bar D$ selects $K=4$).
    \item The $\bar D$+ward partition at $K=4$ is a strict refinement of the in-model $K=2$ PSM partition: each in-model cluster splits cleanly into two, with $14/20/20/14$ sizes.
    \item Adjusted Rand index between in-model $\hat c$ and post-hoc $\bar D$+ward $K=4$: $0.51$.
\end{itemize}
The agreement between (i) the in-model PSM and (ii) post-hoc $\bar D$ confirms that the visible cluster structure is genuine; the disagreement at finer resolutions confirms that the lock-in described in Section~\ref{sec:symptom} is what limits the in-model output, not a real lack of structure.
This is the empirical justification for Fix~A.

\section{Future direction: two-way coupling}\label{sec:future}

The present V8 sacrifices full posterior coherence for ease of implementation. The natural next step replaces the one-way coupling with a fully joint posterior in which the clustering layer feeds back into the LSIRM updates. Two equivalent specifications:

\paragraph{(i) Likelihood feedback.}
Include $N_n(x_j;\delta+\Lambda\eta_j,\Psi)$ in the MH acceptance ratios for $a_i$ and $b_j$:
\[
\alpha(b_j,b_j^\ast)
=
\min\left\{1,\;
\frac{p(Y_{\cdot j}\mid b_j^\ast)\,p(b_j^\ast)\prod_i N(\log\|a_i-b_j^\ast\|;\delta_i+\lambda_i^\top\eta_j,\psi_i)}
{p(Y_{\cdot j}\mid b_j)\,p(b_j)\prod_i N(\log\|a_i-b_j\|;\delta_i+\lambda_i^\top\eta_j,\psi_i)}
\right\}.
\]

\paragraph{(ii) Cluster-conditional prior on $b_j$.}
$b_j\mid c_j=l\sim N_d(\mu_l^{\mathrm{LSIRM}},\Sigma_l^{\mathrm{LSIRM}})$, with corresponding hyperpriors. Items in the same cluster are then explicitly pulled together in the LSIRM latent space, so cluster signal enters the model upstream of the FA.

Either change makes the FMC influence the LSIRM and is expected to alleviate the $\Lambda\to 0$ failure mode by tying $\eta_j$ to a structure that is forced into the LSIRM positions themselves rather than discovered from a noisy log-distance representation.

\section{Implementation summary}

\begin{itemize}[leftmargin=1.5em]
    \item \texttt{my\_LSIRM\_FMC\_v8.cpp}: C++ Gibbs/MH implementation. LSIRM block is byte-identical to V6.
    \item \texttt{my\_LSIRM\_FMC\_cpp\_v8.R}: R wrapper, prior defaults, storage toggles ($\eta$ saved by default; $\Lambda,\delta,\psi$ posterior means only).
    \item \texttt{MIDUS\_5layered\_result\_v8.R}: end-to-end run script for the MIDUS analysis. Uses the PCA initialization of Section~\ref{sec:init} and the hyperparameters of Section~\ref{sec:hyper-rationale}, including the tightened $\psi$ prior.
    \item \texttt{factor\_analysis\_v8/04\_fmc\_factor\_clustering.R}: post-hoc clustering script implementing Section~\ref{sec:posthoc}.
\end{itemize}

\end{document}
